% !TeX root = ../../Book1.tex
\chapter*{后记}
\addcontentsline{toc}{chapter}{后记}

写到这里，第一卷的任务告一段落。从可测集合到低正则函数和不规则边界，所处理的对象逐渐扩大。回看这些证明，反复出现的几种办法也渐渐清楚起来：用简单对象逼近复杂对象，先取得可以控制的估计，再考察这些估计能否保留到极限中。同一种方法换了对象，往往就需要重新检查它成立的条件。

\ 

本卷有意保留测度延拓、积分构造和若干逼近定理的完整证明，因为这些地方能够切实训练良定义、可数运算和例外集的处理。无穷测度、无界函数与不完备空间带来的困难，也尽量分别说明，不一概交给一句“标准处理”。至于细拓扑及流理论中的若干深层结果，本书采用证明概要与明确的文献出处，让读者知道还需要补足什么。这种取舍是为了让已经展开的工具连成一条可以跟随的论证路线，也为更深入的学习留下位置。

\ 

我尤其希望读者从证明中带走一种耐心。许多错误并不发生在最长的计算里，而出在一句未经检查的极限交换、一个没有说明可测性的集合，或一次把有限测度结论直接搬到无穷测度空间的尝试。书中反复展开截断、局部化和对角选择，多少显得笨重，却能让这些缝隙露出来。等方法真正熟练以后，读者自然可以写得更短。就这些初学时容易略去的检查而言，短是结果，不是起点。

\ 

读完本卷，我希望读者能够从生成数据构造测度和积分，分辨几种常用收敛，运用 \(L^p\) 空间与弱收敛的基本工具；面对局部平均、低正则函数或不规则集合时，也知道应当尝试哪一种覆盖、逼近或紧性方法。检验这些收获，不妨合上书，重建一条证明链。能说清每项假设在哪里发挥作用，才更容易判断一个结论能否用于眼前的问题。

\ 

复习时可以从三条路线中任选一条。第一条从前测度走到 Lebesgue 积分与收敛定理；第二条从覆盖引理走到测度微分、面积公式和可整流性；第三条从分布导数走到 Sobolev、BV 与有限周长集合。卡住时先查对象所在的空间和所用的收敛，再查估计是否一致、例外集是否仍可控。章末习题也值得隔一段时间脱离原解重做；那些读答案时顺畅无碍、动笔时却接不上的地方，往往最能说明接下来该补什么。

\ 

继续深入测度的结构与反例，可读 Bogachev 的《Measure Theory》\cite{Bogachev2007Measure}；几何测度方向可沿 Federer 的《Geometric Measure Theory》\cite{Federer1996Geometric}和 Mattila 的《Geometry of Sets and Measures in Euclidean Spaces: Fractals and Rectifiability》\cite{Mattila1995Geometry}继续。若更关心 Sobolev 与 BV，Leoni 的《A First Course in Sobolev Spaces》\cite{Leoni2017First}可作为进一步学习的教材，Ambrosio、Fusco 与 Pallara 的《Functions of Bounded Variation and Free Discontinuity Problems》\cite{Ambrosio2000Functions}则通向更专门的理论。这些书不必同时铺开，也不必一律从头读到尾。选定一条路线，持续读完与当前问题相关的部分，再逐步扩大范围，更容易积累起可以调用的工具。

\ 

若问题来自方程、频率或振荡，可以接着读第二卷；若兴趣落在随机现象，可从本卷的测度与积分转入概率论。研究界面、奇异集合和变分问题的读者，则可以沿几何测度与 BV 继续。进入专著或论文时，不妨随手记下所用结果之间的依赖：哪些已经掌握，哪些只在有限测度或光滑对象上成立，哪些步骤还需补证。遇到暂时不懂的地方，先把缺口认清，也是一种进展。

\ 

一本基础书能做的事终究有限。它可以整理语言，留下若干可靠的工具，也可以让一些常见的错误更早暴露；真正的熟练仍要在具体问题中取得。读到这里，下一页不必再属于这本书了。

\ 

赠读者：

\begin{quote}
\noindent
世上无难事，只要肯登攀。\\
\hfill——摘自毛泽东《水调歌头·重上井冈山》
\end{quote}
